\hypertarget{appendix-736a2345de4a}{}\label{referencestart:appendix-736a2345de4a}
\section{Mixed signedness}\label{appendix-736a2345de4a-mixed-signedness}

\hypertarget{appendix-736a2345de4a-mixed-signedness}{}

\textbf{Curriculum depth:} Introduced · \textbf{First seen:} C002 · \textbf{Latest development:} C002

\begin{allmightroledefinition}

\subsection{Definition and mechanism}\label{appendix-736a2345de4a-definition-and-mechanism}

\hypertarget{appendix-736a2345de4a-definition-and-mechanism}{}

A negative signed operand can convert to a large unsigned value; choose semantic types consistently or use \BookInlineCode{std::cmp\_*}, \BookInlineCode{std::in\_range}, or \BookInlineCode{std::ssize} where appropriate.

\end{allmightroledefinition}

int −1 compared with unsigned int 1 converts to a large unsigned value on the usual host, so the comparison is false. std::cmp\_less preserves mathematical ordering. A valid index still needs a lower bound: \BookInlineCode{index\textgreater=0\ \&\&\ std::cmp\_less(index,size)}.

\subsection{Prerequisites}\label{appendix-736a2345de4a-prerequisites}

\hypertarget{appendix-736a2345de4a-prerequisites}{}

\begin{itemize}
\tightlist
\item
  \hyperref[appendix-5e2d663adf7f]{Usual arithmetic conversions} --- After promotions, equal-signedness operands use the higher-ranked type.
\end{itemize}

\subsection{Used by}\label{appendix-736a2345de4a-used-by}

\hypertarget{appendix-736a2345de4a-used-by}{}

Connect later chunks here when they use this concept.

\subsection{Recall and next depth}\label{appendix-736a2345de4a-recall-and-next-depth}

\hypertarget{appendix-736a2345de4a-recall-and-next-depth}{}

Explain mixed signedness without the label, work the example above, then state the assumption whose removal would change the result.

\textbf{Pending:} API design, diagnostics, and range-safe indexing.

\subsection{Source and navigation}\label{appendix-736a2345de4a-source-and-navigation}

\hypertarget{appendix-736a2345de4a-source-and-navigation}{}

\hyperref[01-foundations-i-ch001-arithmetic-becomes-a-policy]{C002 teaching chapter} · \href{../01-Bits-Programs-and-Reliable-Systems-P03/main.pdf\#nameddest=appendix-fd090d982a15}{Part definitions and mechanisms (Bits, Programs, and Reliable Systems, Part 3, Reference notes: Foundations)} · \href{../07-Vision-XR-and-Playable-Worlds-P01/main.pdf\#nameddest=appendix-4ad47bdd269a}{Dependency Map (Vision, XR, and Playable Worlds, Part 1, Reference notes: Dependency Map)}
